\documentclass[10pt,a4paper]{amsart}
\usepackage[margin=19mm]{geometry}
\usepackage{amsmath,amssymb,mathtools}
\usepackage[scr=rsfs,cal=euler]{mathalfa}
\usepackage{xfrac}
\usepackage[unicode,hidelinks,bookmarks=false]{hyperref}
\newcommand{\ie}{i.e.\ }
\newcommand{\cf}{cf.\ }
\newcommand{\resp}{respectively\ }
\newcommand{\Subsection}{\subsection}
\providecommand{\nobreakdash}{\nobreak\hbox{-}}
\setlength{\emergencystretch}{2em}
\multlinegap0pt
\usepackage{bm}
\usepackage{bbm}
\usepackage{dsfont}
\usepackage{xspace}
\usepackage{cancel}
\usepackage{mathtools}
\usepackage{amsthm}
\makeatletter
\@ifundefined{theorem}{\newtheorem{theorem}{Theorem}}{}
\@ifundefined{definition}{\newtheorem{definition}[theorem]{Definition}}{}
\@ifundefined{lemma}{\newtheorem{lemma}[theorem]{Lemma}}{}
\makeatother
\title[Existence and Nonexistence Breaking Results For a Weighted E]{Existence and Nonexistence Breaking Results For a Weighted Elliptic Problem in Half-Space}
\author{J. M. Do \'O, R. F. Freire, J. Giacomoni, E. S. Medeiros}
\date{Source: arXiv:2510.05999v1 (CC BY 4.0)}
\begin{document}
\maketitle

\noindent\emph{Source and licence.} Existence and Nonexistence Breaking Results For a Weighted Elliptic Problem in Half-Space. Version arXiv:2510.05999v1 (CC BY 4.0), distributed under the Creative Commons Attribution 4.0 International licence, as recorded in the arXiv licence metadata for this exact version and shown on the abstract page https://arxiv.org/abs/2510.05999v1. Full rights notice: licenses/arxiv-2510.05999-CC-BY-4.0.txt.\par\smallskip

\noindent\textit{Editorial scope.} Counted unit: the Lemma whose source label is \texttt{Linfty 2} in the article (source0.tex lines 1015--1022, source label \texttt{Linfty 2}), delivered together with the article's own complete original proof of it (source0.tex lines 1024--1148, one \texttt{proof} environment of 125 lines). No mathematics is written, repaired or re-derived: statement and proof are the article's own text line for line. Required context retained uncounted: PG (lines 112--122); weak solution (lines 348--359); Linfty (lines 916--923). Every definition, display and label of those lines is retained unchanged. Omitted from this package and not redistributed: 1--111, 123--347, 360--915, 924--1014, 1023--1023, 1149--1847. Nothing inside a retained excerpt is omitted or altered: every retained body is delivered exactly as the article's lines are written, so no omission receipt of any kind accompanies this package. Citations: the retained excerpts carry no citation command at all, so no bibliography entry of the article is transcribed and no reference is redistributed. Reference adaptation: none was needed and none was made; every reference inside the delivered unit resolves inside it, so no unresolved reference or citation is produced. Printed slips kept literal: no printed word, symbol, hypothesis or number of the counted statement or of its proof was altered. Layout and class: the article's own class is \texttt{amsart} with packages including xcolor, amsfonts, color, latexsym, amssymb, amsmath, inputenc, dsfont, enumitem, float, graphicx, tikz, tkz-euclide, geometry; the wrapper is the lane's pinned \texttt{amsart} editorial wrapper at 10pt on A4, which restates the theorem-like environments the retained text needs and adds the article's own macro definitions that the retained text needs (none), with extra wrapper packages (bm, bbm, dsfont, xspace, cancel, mathtools). The delivered numbering is the unit's own. Assets: no retained span contains a figure environment, a graphics-inclusion command, a PDF-page inclusion, a table or a TikZ picture; no image file is copied into this package, the package declares no asset dependency and no image file is redistributed with it. Licence: version arXiv:2510.05999v1 of this article is distributed under the Creative Commons Attribution 4.0 International licence, as recorded in the arXiv licence metadata for this exact version and shown on the abstract page https://arxiv.org/abs/2510.05999v1. A full rights notice is delivered at licenses/arxiv-2510.05999-CC-BY-4.0.txt. Characters: every retained excerpt is pure ASCII, so the delivered engine, XeLaTeX with its default Latin Modern text font, reports no missing character.
  \begin{equation}\label{PG}
		\left\{
		\begin{aligned}
			-\mathrm{div}(\rho(x_N) \nabla u) &=a|u|^{p-2}u &\mbox{in }&\ \mathds{R}^N_+,
			\vspace{0.2cm}\\
			-\frac{\partial u}{\partial x_N}&=b|u|^{q-2}u&\mbox{on }&\
			\mathds{R}^{N-1},
		\end{aligned}
		\right. 
		\tag{$\mathcal{P}_0$}
	\end{equation}

\begin{definition} By a weak solution of \eqref{PG} we mean a function  
$ u \in \mathcal{D}^{1,2}_\rho(\mathds{R}^N_+)$
such that  
\begin{equation}\label{weak solution}
   \int_{\mathds{R}^{N}_+} \rho(x_N)\,\nabla u \cdot \nabla \varphi \, \mathrm{d}x
   = a \int_{\mathds{R}^{N}_+} |u|^{p-2}u \varphi \, \mathrm{d}x  
   + b \int_{\mathds{R}^{N-1}} |u|^{q-2}u \varphi \, \mathrm{d}x^\prime,
\end{equation}
for every testing function $\varphi \in C_\delta^\infty(\mathds{R}^N_+)$.  
For simplicity, and without loss of generality, we assume throughout this work that $\rho(0)=1$.

\end{definition}

\begin{lemma}\label{Linfty}
Assume that condition $(\rho_0)$ holds. Then, in each of following assertions:
\begin{itemize}
    \item[$(i)$]  $u$ is a nonnegative weak solution in $E_q(\mathds{R}^{N-1})$, with $\gamma\geq 2$, $a > 0$, $b \leq 0$, $p \in [2,2^\ast)$, and $q \in (1,\infty)$;
    \item[$(ii)$]  $u$ is a nonnegative weak solution in $E_p(\mathds{R}^N_+)$, with $\gamma>1$, $a \leq 0$, $b > 0$, $p \in (1,\infty)$, and $q \in [2, 2_\ast)$,
\end{itemize}
we have $u \in L^\infty(\mathds{R}^N_+) \cap L^\infty(\mathds{R}^{N-1})$.
\end{lemma}

\begin{lemma}\label{Linfty 2}
Assume condition $(\rho_0)$ holds with $\gamma\geq 0$. Then, in each of the following cases:
\begin{itemize}
\item[$(i)$] $u$ is a nonnegative weak solution in $E_q(\mathds{R}^{N-1})$, with $a > 0$, $b \leq 0$, $p=2^\ast$ and $q\in(1,\infty)$;
\item[$(ii)$] $u$ is a nonnegative weak solution in $E_p(\mathds{R}^N_+)$, with $a \leq 0$, and $b > 0$, $p\in(1,\infty)$ and $q=2_\ast$,
\end{itemize}
it follows that $u \in L^\infty(\mathds{R}^N_+) \cap L^\infty(\mathds{R}^{N-1})$.
\end{lemma}

\begin{proof} Let us prove the first item. By the classical Sobolev inequality together with assumption $(\rho_0)$, we also have $E_q(\mathds{R}^{N-1})\hookrightarrow \mathcal{D}^{1,2}_\rho(\mathds{R}^N_+)\hookrightarrow L^{2^\ast}(\mathds{R}^N_+)$. Then, arguing by density the definition of weak solution \eqref{weak solution} holds for testing functions in $E_q(\mathds{R}^{N-1})$. Let $\varphi:[0,\infty)\longrightarrow \mathds{R}$ defined by 
\begin{equation*}
    \varphi(t)= \varphi_{m,k}(t)=\left\lbrace
		\begin{aligned}
			&t^k,\quad&\mbox{for }&\,t\leq m\\
		&km^{k-1}(t-m)+m^k,\quad &\mbox{for}&\,t\geq m,
		\end{aligned}
		\right. 
\end{equation*}
for $k>1$ and $m>0$.
Then we have
\begin{equation*}
    \varphi^\prime(t)=\left\lbrace
		\begin{aligned}
			&kt^{k-1},\quad&\mbox{for }&\,t\leq m\\
		&km^{k-1},\quad &\mbox{for}&\,t\geq m.
		\end{aligned}
		\right. 
\end{equation*}
Take $\varphi(u)\varphi'(u)$ as a testing function in \eqref{weak solution}, then, given that $\varphi(u)\varphi'(u)\geq0$ and $b\leq 0$, we obtain 
\begin{align}\label{first inequality}
   \nonumber \int_{\mathds{R}^{N}_+}\rho(x_N)\nabla u \nabla (\varphi(u)\varphi^\prime(u)) \, \mathrm{d} x&=a\int_{\mathds{R}^{N}_+}u^{2^\ast-1}\varphi(u)\varphi^\prime(u) \mathrm{d}x  +b\int_{\mathds{R}^{N-1}}u^{q-1} \varphi(u)\varphi^\prime(u) \, \mathrm{d} x\\
    &\leq a\int_{\mathds{R}^{N}_+}u^{2^\ast-1} \varphi(u)\varphi^\prime(u) \, \mathrm{d} x.
\end{align}
Now, observe that
\begin{equation*}
    \nabla \varphi(u)=\left\lbrace
		\begin{aligned}
			&ku^{k-1}\nabla u,\quad&\mbox{for }&\,u\leq m\\
		&km^{k-1}\nabla u,\quad &\mbox{for}&\,u\geq m
		\end{aligned}
		\right. 
\end{equation*}
and
\begin{equation*}
    \nabla (\varphi(u)\varphi^\prime(u))=\left\lbrace
		\begin{aligned}
			&k(2k-1)u^{(2k-1)}\nabla u,\quad&\mbox{for }&\,u\leq m\\
		&k^2 m^{2(k-1)}\nabla u,\quad &\mbox{for }&\,u\geq m.
		\end{aligned}
		\right. 
\end{equation*}
Then $\nabla u \nabla (\varphi(u)\varphi^\prime(u))\geq |\nabla \varphi(u)|^2$ and
\begin{equation}\label{LHS}
    \int_{\mathds{R}^N_+}\rho(x_N)\nabla u\nabla(\varphi(u)\varphi^\prime(u))\mathrm{d}x\geq \int_{\mathds{R}^N_+}\rho(x_N)|\nabla \varphi(u)|^2\mathrm{d}x\geq C\left(\int_{\mathds{R}^{N}_+}\varphi(u)^{2^\ast}\mathrm{d}x\right)^{2/2^\ast}.
\end{equation}
On the other hand, since $u\varphi^\prime(u)\leq k\varphi(u)$, we have
\begin{equation*}
    \int_{\mathds{R}^{N}_+}u^{2^\ast-1}\varphi(u)\varphi^\prime(u)\mathrm{d}x\leq k\int_{\mathds{R}^{N}_+}\varphi(u)^2u^{2^\ast-2}\mathrm{d}x.
\end{equation*}
Thus, applying the inequality above and \eqref{LHS} in \eqref{first inequality} we obtain
\begin{equation}\label{varphi inequality}
    \left(\int_{\mathds{R}^N_+}\varphi(u)^{2^\ast}\mathrm{d}x\right)^{2/2^\ast}\leq C_1k\int_{\mathds{R}^N_+}\varphi(u)^2u^{2^\ast-2}\mathrm{d}x, \quad\forall k>1.
\end{equation}

\textbf{Claim:} Let $k_1$ such that $2k_1=2^\ast$. Then $u\in L^{k_12^\ast}(\mathds{R}^N_+)$.

Let $R>0$ to be determined later. By Holder's inequality we obtain
\begin{align*}
    \int_{\mathds{R}^N_+}\varphi(u)^2u^{2^\ast-2}\mathrm{d}x=&\int_{\{u\leq R\}}\varphi(u)u^{2^\ast-2}\mathrm{d}x+\int_{\{u>R\}}\varphi(u)^2u^{2^\ast-2}\mathrm{d}x\\
    \leq& R^{2^\ast-2}\int_{\mathds{R}^N_+}\varphi(u)^2\mathrm{d}x+\left(\int_{\mathds{R}^N_+}\varphi(u)^{2^\ast}\mathrm{d}x\right)^{2/2^\ast}\left(\int_{\{u>R\}}u^{2^\ast}\mathrm{d}x\right)^{(2^\ast-2)/2^\ast}.
\end{align*}
By dominated convergence Theorem we can choose $R>0$ such that 
\begin{equation*}
    \left(\int_{\{u>R\}}u^{2^\ast}\mathrm{d}x\right)^{(2^\ast-2)/2^\ast} \leq \frac{1}{2C_1k_1}.
\end{equation*}
Thus,
\begin{equation*}
    \int_{\mathds{R}^N_+}\varphi(u)^2u^{2^\ast-2}\mathrm{d}x\leq R^{2^\ast-2}\int_{\mathds{R}^N_+}\varphi(u)^2\mathrm{d}x+\frac{1}{2C_1k_1}\left(\int_{\mathds{R}^N_+}\varphi(u)^{2^\ast}\mathrm{d}x\right)^{2/2^\ast}.
\end{equation*}
Applying this inequality in \eqref{varphi inequality} we obtain
\begin{equation*}
    \left(\int_{\mathds{R}^N_+}\varphi(u)^{2^\ast}\mathrm{d}x\right)^{2/2^\ast}\leq C_1k_1 R^{2^\ast-2}\int_{\mathds{R}^N_+}\varphi(u)^2\mathrm{d}x.
\end{equation*}
Hence, given that $\varphi_{m,k}(u)\leq u^{k}$, letting $m\longrightarrow\infty$, by dominated convergence Theorem we obtain
\begin{equation*}
     \left(\int_{\mathds{R}^N_+}u^{k_12^\ast}\mathrm{d}x\right)^{2/2^\ast} \leq C_1k_1R^{2^\ast-2}\int_{\mathds{R}^N_+}u^{2^\ast}\mathrm{d}x<\infty
\end{equation*}
and we conclude the claim.

Applying dominated convergence Theorem again in \eqref{varphi inequality} we obtain
\begin{equation*}
    \left(\int_{\mathds{R}^N_+}u^{k2^\ast}\mathrm{d}x\right)^{2/2^\ast}\leq C_1 k \int_{\mathds{R}^N_+}u^{2k+2^\ast-2}\mathrm{d}x,
\end{equation*}
and consequently
\begin{equation}\label{bootstrap}
    \left(\int_{\mathds{R}^N_+}u^{k2^\ast}\mathrm{d}x\right)^{1/2^\ast(k-1)}\leq (C_1k)^{1/2(k-1)} \left(\int_{\mathds{R}^N_+} u^{2k
    +2^\ast-2}\mathrm{d}x\right)^{1/2(k-1)},\quad \forall k>1.
\end{equation}
Define $2k_{n}+2^\ast-2=2^\ast k_{n-1}$, then
\begin{equation*}
k_n-1=\frac{2^\ast}{2}(k_{n-1}-1)=\left(\frac{2^\ast}{2}\right)^{n-1}(k_1-1)
\end{equation*}
and $k_n\longrightarrow \infty$. Therefore, we have
\begin{equation*}
    \left(\int_{\mathds{R}^N_+}u^{k_n2^\ast}\mathrm{d}x\right)^{1/2^\ast(k_n-1)}\leq (C_1k_n)^{1/2(k_n-1)} \left(\int_{\mathds{R}^N_+} u^{k_{n-1}2^\ast}\mathrm{d}x\right)^{1/2^\ast(k_{n-1}-1)},
\end{equation*}
and applying recursively the inequality \eqref{bootstrap} we obtain
\begin{equation*}
    \left(\int_{\mathds{R}^N_+}u^{k_n2^\ast}\mathrm{d}x\right)^{1/2^\ast(k_n-1)}\leq \prod_{i=2}^n(C_1k_i)^{1/2(k_i-1)} \left(\int_{\mathds{R}^N_+} u^{k_{1}2^\ast}\mathrm{d}x\right)^{1/2^\ast(k_{1}-1)}.
\end{equation*}
Write $k_i^{1/2(k_i-1)}=\left[k_i^{1/\sqrt{k_i+1}}\right]^{1/2\sqrt{k_i+1}}$ and note that $\displaystyle\lim_{k\longrightarrow \infty}k^{1/\sqrt{k-1}}=1$, then there exists $C_2>1$ such that
\begin{equation*}
  \left(\int_{\mathds{R}^N_+}u^{k_n2^\ast}\mathrm{d}x\right)^{1/2^\ast(k_n-1)}\leq C_1^{\sum_{i=2}^n{\frac{1}{2(k_i-1)}}} C_2^{\sum_{i=2}^n{\frac{1}{2\sqrt{k_i-1}}}} \left(\int_{\mathds{R}^N_+} u^{k_{1}2^\ast}\mathrm{d}x\right)^{1/2^\ast(k_{1}-1)}.
  \end{equation*}
Since
\begin{equation*}
    \sum_{n=2}^\infty \frac{1}{2(k_n-1)}=\frac{1}{2(k_1-1)} \sum_{n=2}^\infty \left(\frac{2}{ 2^\ast}\right)^{n-1}<\infty  
\end{equation*}
and
\begin{equation*}
    \sum_{n=2}^\infty \frac{1}{2\sqrt{k_n-1}}=\frac{1}{2\sqrt{k_1-1}}\sum_{n=2}^\infty \left(\sqrt{\frac{2}{ 2^\ast}}\right)^{n-1}<\infty,
\end{equation*}
we obtain
\begin{equation*}
    \left(\int_{\mathds{R}^N_+}u^{k_n2^\ast}\mathrm{d}x\right)^{1/2^\ast(k_n-1)}\leq C_0\|u\|_{k_12^\ast,\mathds{R}^N_+}^{\frac{k_1}{k_1-1}}.
\end{equation*}
Therefore, letting $n\longrightarrow \infty$ we obtain $u\in L^{\infty}(\mathds{R}^N_+)$ with $\|u\|_{\infty, \mathds{R}^N_+}\leq C_0 \|u\|_{k_12^\ast,\mathds{R}^N_+}^{\frac{k_1}{k_1-1}}$. Finally, arguing as in Lemma \ref{Linfty}, we obtain $u\in L^\infty(\mathds{R}^{N-1})$ and conclude the first item.

For the second item, a similar argument shows that $u \in L^\infty(\mathds{R}^{N-1})$, with the bound
\begin{equation*}
\|u\|_{\infty,\mathds{R}^{N-1}}\leq C_0 \|u\|_{k_12_\ast,\mathds{R}^{N-1}}^{\frac{k_1}{k_1-1}}.
\end{equation*}
Finally, following the argument at the end of the second item in Lemma~\ref{Linfty}, we conclude that $u\in L^\infty(\mathds{R}^N_+)$ and complete case $(ii)$. 
\end{proof}

\end{document}
